\chapter{Novelty Gap Analysis}
\label{ch:novelty}

A journal contribution has to be new against the work that exists on the day
of submission, not on the day the project was proposed. This chapter maps 64
papers from 2024 to September 2026 that are close to the project. It
identifies which directions are saturated and states precisely the one that
remains open.

\section{A saturated design space}
\label{sec:saturated}

Figure~\ref{fig:timeline} places the papers on a timeline in four lanes.
\begin{itemize}
  \item The first lane, \emph{Mamba inside YOLO or for small objects},
        holds 24 papers. Three of them appeared in September 2026 alone:
        ScopeMamba-YOLO, YOLO12-MambaScan and HDMamba-YOLO~\citep{scopemamba2026,yolo12mambascan2026,hdmambayolo2026}.
        There is also a CVPR 2026 paper~\citep{akcmamba2026} and a steady
        flow of journal articles.
  \item The second lane, \emph{Mamba for visible--thermal fusion}, holds 17
        papers. They appeared in ICCV, IEEE TMM, IEEE TCSVT and two volumes
        of Information Fusion.
  \item The third lane covers state-space models used \emph{across time} for
        any task. It is dominated by tracking and by event cameras.
  \item The fourth lane covers \emph{visible--thermal video detection}
        without state-space models. It uses windows, graphs and hypergraphs.
\end{itemize}

The design implied by a literal reading of the project title is a YOLO
detector with Mamba blocks fusing visible and thermal features on still
images. It sits at the intersection of the first two lanes. Fusion-Mamba
describes itself as ``a dual-stream feature extraction network and three
Fusion-Mamba blocks, with the same neck and head as
YOLOv8''~\citep{dong2025fusionmamba}. MambaRefine-YOLO, RemoteDet-Mamba,
WaveMamba and COMO occupy the same point~\citep{mambarefine2025,remotedetmamba2024,wavemamba2025,como2024}.
A reviewer at IEEE TCSVT or Pattern Recognition would read a further
instance as incremental.

\begin{figure}[htbp]
\centering
\begin{tikzpicture}[
    pd/.style={circle, draw=none, inner sep=0pt, minimum size=2.0mm},
    pl/.style={font=\scriptsize\itshape, text=muted, inner sep=1pt},
  ]
  % 1 month = \u cm, month index m = 12*(year-2024) + (month-1)
  \def\u{0.285}
  \def\yE{5.6}\def\yA{4.2}\def\yB{2.8}\def\yC{1.4}\def\yD{0}
  \newcommand{\pt}[3]{\node[pd, fill=#3] at ({#1*\u},{#2}) {};}

  % lane labels
  \node[grp, anchor=east, align=right] at (-2.05,\yE) {this work};
  \node[grp, anchor=east, align=right] at (-2.05,\yA) {Mamba in YOLO\\ or small objects};
  \node[grp, anchor=east, align=right] at (-2.05,\yB) {Mamba for\\ RGB--T fusion};
  \node[grp, anchor=east, align=right] at (-2.05,\yC) {SSM across time,\\ any task};
  \node[grp, anchor=east, align=right] at (-2.05,\yD) {RGB--T video,\\ no SSM};

  % lane A: Mamba in YOLO / small objects
  \foreach \m/\o in {4/0, 4/0.24, 5/0, 7/0, 8/0, 17/0, 18/0, 21/0, 22/0, 24/0, 27/0,
                     29/0, 29/0.24, 30/0, 30/0.24, 30/0.48, 32/0, 32/0.24, 32/0.48}
    \pt{\m}{\yA+\o}{cRGB!75};
  \node[pl, anchor=north] at ({5*\u},{\yA-0.15}) {Mamba YOLO};
  \node[pl, anchor=north] at ({24*\u},{\yA-0.15}) {UAVDet};
  \node[pl, anchor=north] at ({30*\u},{\yA-0.15}) {MambaPSA};
  \node[pl, anchor=south] at ({32*\u},{\yA+0.6}) {three in Sept.\ 2026};

  % lane B: Mamba for RGB-T fusion
  \foreach \m/\o in {3/0, 3/0.24, 6/0, 7/0, 9/0, 11/0, 18/0, 18/0.24, 18/0.48, 20/0,
                     21/0, 22/0, 23/0, 24/0, 26/0}
    \pt{\m}{\yB+\o}{cFuse!80};
  \node[pl, anchor=north] at ({3*\u},{\yB-0.15}) {Fusion-Mamba, CFMW};
  \node[pl, anchor=north] at ({11*\u},{\yB-0.15}) {COMO};
  \node[pl, anchor=south] at ({18*\u},{\yB+0.6}) {WaveMamba, MS2Fusion};
  \node[pl, anchor=north] at ({22.5*\u},{\yB-0.15}) {MambaRefine};

  % lane C: SSM across time
  \foreach \m/\o in {1/0, 6/0, 7/0, 7/0.24, 8/0, 9/0, 9/0.24, 10/0, 12/0, 21/0, 29/0, 31/0}
    \pt{\m}{\yC+\o}{cTime!80};
  \node[pl, anchor=north] at ({1*\u},{\yC-0.15}) {Zubi\'c et al.};
  \node[pl, anchor=south] at ({9*\u},{\yC+0.37}) {SambaMOTR};
  \node[pl, anchor=north] at ({29*\u},{\yC-0.15}) {TMambaDet};

  % lane D: RGB-T video detection without SSM
  \foreach \m in {-5, 9, 15, 24, 30}
    \pt{\m}{\yD}{cIR!80};
  \node[pl, anchor=north] at ({-5*\u},{\yD-0.15}) {EINet};
  \node[pl, anchor=north] at ({9*\u},{\yD-0.15}) {MSENet};
  \node[pl, anchor=north] at ({15*\u},{\yD-0.15}) {MSGNet};
  \node[pl, anchor=north] at ({24*\u},{\yD-0.15}) {Strip-Fusion};
  \node[pl, anchor=north] at ({30*\u},{\yD-0.15}) {DCHNet};

  % this work: empty lane with a target marker
  \node[star, star points=5, star point ratio=2.3, fill=cAlert, draw=none,
        inner sep=0pt, minimum size=3.6mm] at ({33.6*\u},\yE) {};
  \node[pl, anchor=east, text=cAlert] at ({33.0*\u},\yE)
       {causal SSM memory, RGB--T streams, tiny objects: no prior work found};

  % time axis: ticks and year labels only
  \foreach \m/\lab in {-6/{Jul 2023}, 0/{Jan 2024}, 12/{Jan 2025}, 24/{Jan 2026}, 33/{Oct 2026}}{
    \draw[draw=cNeutral, line width=0.5pt] ({\m*\u},-0.85) -- ({\m*\u},-0.7);
    \node[font=\scriptsize, text=muted, anchor=north] at ({\m*\u},-0.88) {\lab};
  }
\end{tikzpicture}

\vspace{2mm}
\caption[Timeline of closely related work, 2023 to 2026]{Timeline of
closely related work from July 2023 to October 2026, one dot per paper,
stacked when several share a month. Papers without a known month of
publication are omitted. Static Mamba-in-YOLO and Mamba fusion papers
accumulate quickly. Temporal state-space work concentrates on tracking and
event cameras. Visible--thermal video detection has advanced without
state-space models. No work occupies the top lane.}
\label{fig:timeline}
\end{figure}


\section{Five candidate directions}
\label{sec:directions}

Table~\ref{tab:directions} assesses five directions that would extend the
static design. Each is judged by the closest prior art found and by what
remains open.

\begin{table}[htbp]
\centering
\caption[Candidate novelty directions and their prior art]{Candidate novelty
directions and their prior art. Risk is the likelihood that a reviewer finds
the contribution already published.}
\label{tab:directions}
\small
\begin{tabularx}{\textwidth}{@{}P{2.6cm}P{3.6cm}LP{1.35cm}@{}}
\hd{Direction} & \hd{Closest prior art} & \hd{What remains open} & \hd{Risk} \\ \gap
(a) causal state carried across frames & MambaST~\citep{mambast2024}; Zubi\'c et al.~\citep{zubic2024ssmevent}; TMambaDet~\citep{tmambadet2026}; StreamYOLO~\citep{yang2022streamyolo} & No causal, constant-memory, unbounded-length visible--thermal detector. No state-space method on RGB--thermal video detection benchmarks. & low to medium \\ \sgap
(b) misalignment and missing modality & COMO~\citep{como2024}; contextually guided SSM fusion~\citep{cgssm2025}; MSDF-Mamba~\citep{msdfmamba2025}; SCDT for tracking~\citep{scdt2026} & Temporal missing modality in detection, such as bursty thermal dropout and desynchronisation, has no benchmark protocol. & medium; low if combined with (a) \\ \sgap
(c) size-aware scan order & QuadMamba~\citep{xie2024quadmamba}; DAMamba~\citep{damamba2025}; ObjM~\citep{objm2025}; ScopeMamba~\citep{scopemamba2026} & Only a temporal twist, ordering tokens by objects predicted in the previous frame. & high \\ \sgap
(d) joint detection and tracking with shared state & SambaMOTR~\citep{segu2025sambamotr}; MambaTrack~\citep{mambatrack2024}; TrackSSM~\citep{trackssm2024} & A shared state for dense YOLO-style detection and identity in RGB--thermal tiny-object tracking. & medium; large scope \\ \sgap
(e) edge deployment of SSM detectors & MambaNeXt-YOLO~\citep{mambanextyolo2025}; LCMamNet~\citep{lcmamnet2026}; DLRMamba~\citep{dlrmamba2026} & No native TensorRT operator for two-dimensional selective scan was found; accuracy loss after export is rarely reported. & medium to high alone \\
\end{tabularx}
\end{table}

Direction (c) is crowded and should be avoided as a main claim. Directions
(d) and (e) are valuable but either too large for one paper or too narrow
for a top journal. Directions (a) and (b) are complementary. Combined, they
define a problem no published work addresses.

\section{The open problem}
\label{sec:openproblem}

Stated precisely, the open problem is this:

\begin{quote}
\itshape
Detect very small objects in a continuous stream of visible and thermal
frames, using only past frames, at constant memory and per-frame cost, while
remaining accurate when one sensor drops out, drifts out of alignment or
delivers frames at irregular intervals.
\end{quote}

Figure~\ref{fig:positioning} places the closest work on two axes. The first
is how much temporal context a method uses. The second is how demanding its
sensing assumptions are. Every cell to the left of the last column is
occupied. In the last column, causal and unbounded memory exists only for
event cameras and for DETR-style tracking. The two cells in the lower right
are empty. Those cells are the target of \methodS{}.

\begin{figure}[htbp]
\centering
\begin{tikzpicture}[
    cell/.style={base, n, text width=3.35cm, minimum height=2.05cm, font=\scriptsize},
    goal/.style={base, fill=cAlert!11, text width=3.35cm, minimum height=2.05cm, font=\scriptsize},
    rowl/.style={grp, text width=2.4cm, align=right, anchor=east},
    coll/.style={grp, text width=3.4cm},
  ]
  \def\cx{3.75}\def\cy{2.35}
  % column headings
  \node[coll] at (0*\cx,0.05+1.45) {single image};
  \node[coll] at (1*\cx,0.05+1.45) {short window\\[-1pt] of frames};
  \node[coll] at (2*\cx,0.05+1.45) {causal memory,\\[-1pt] unbounded length};
  \node[font=\footnotesize\itshape, text=muted] at (1*\cx,2.35) {temporal context used};

  % row headings
  \node[rowl] at (-1.95,0)       {visible only};
  \node[rowl] at (-1.95,-1*\cy)  {visible and\\ thermal, aligned};
  \node[rowl] at (-1.95,-2*\cy)  {visible and thermal,\\ misaligned or\\ missing frames};

  % row 1: visible only
  \node[cell] at (0*\cx,0) {Mamba YOLO, MambaPSA,\\ UAVDet, HPS-DETR,\\ MambaVision, YOLO26};
  \node[cell] at (1*\cx,0) {TMambaDet (offline),\\ StreamYOLO,\\ LongShortNet,\\ DAMO-StreamNet};
  \node[cell] at (2*\cx,0) {Zubi\'c et al., RVT\\ (event cameras);\\ SambaMOTR\\ (memory per track)};

  % row 2: aligned RGB-T
  \node[cell] at (0*\cx,-1*\cy) {Fusion-Mamba, CFMW,\\ WaveMamba, MS2Fusion,\\ MambaRefine-YOLO};
  \node[cell] at (1*\cx,-1*\cy) {MambaST (3 frames,\\ KAIST only), EINet,\\ Strip-Fusion};
  \node[goal] (g1) at (2*\cx,-1*\cy) {\textbf{empty}\\[2pt] target of \methodS:\\ temporal state memory};

  % row 3: misaligned / missing
  \node[cell] at (0*\cx,-2*\cy) {COMO, contextually\\ guided SSM fusion,\\ MSDF-Mamba, DMM};
  \node[cell] at (1*\cx,-2*\cy) {MSENet, MSGNet,\\ MCDA-Net, DCHNet\\ (alignment-free)};
  \node[goal] (g2) at (2*\cx,-2*\cy) {\textbf{empty}\\[2pt] target of \methodS:\\ reliability-gated $\Delta$};

  \node[font=\footnotesize\itshape, text=muted, rotate=90] at (-4.55,-1*\cy) {sensing assumptions, increasingly demanding};
\end{tikzpicture}

\vspace{2mm}
\caption[Positioning of the closest work by temporal context and sensing]{Positioning
of the closest work. Columns give the temporal context a method uses.
Rows give how demanding its sensing assumptions are. Each cell lists
representative methods. The two highlighted cells contain no published
method that the searches could find. They are the targets of \methodS{}.}
\label{fig:positioning}
\end{figure}


\section{Positioning against the nearest work}
\label{sec:nearest}

Six papers come closest, and the contribution must be defined against each
(Table~\ref{tab:nearest}).
\begin{itemize}
  \item MambaST is the nearest in spirit~\citep{mambast2024}. It fuses
        visible and thermal features across three frames with a stride of
        three. It feeds the last hidden output into the next window and is
        evaluated on KAIST only. Its authors name longer sequences and other
        datasets as future work.
  \item Zubi\'c and colleagues show truly recurrent state-space detection
        with frequency generalisation, on event data~\citep{zubic2024ssmevent}.
  \item TMambaDet is bidirectional and therefore offline~\citep{tmambadet2026}.
  \item SambaMOTR keeps state per track, not per location. It is built on a
        DETR tracker~\citep{segu2025sambamotr}.
  \item COMO handles misalignment on still images only~\citep{como2024}.
  \item RVT is a recurrent detector for event cameras built on convolutional
        LSTM cells, not state-space models~\citep{gehrig2023rvt}.
\end{itemize}

\begin{table}[htbp]
\centering
\caption[Properties of the nearest prior work]{Properties of the nearest
prior work against the proposed \methodS{}.}
\label{tab:nearest}
\small
\begin{tabularx}{\textwidth}{@{}LC{1.25cm}C{1.25cm}C{1.25cm}C{1.25cm}C{1.35cm}C{1.25cm}@{}}
\hd{Property} & \hd{MambaST} & \hd{Zubi\'c} & \hd{TMamba-Det} & \hd{Samba-MOTR} & \hd{COMO} & \hd{\methodS} \\ \gap
detection of objects          & yes & yes & yes & tracks & yes & yes \\
uses only past frames         & yes & yes & no  & yes & n/a & yes \\
memory unbounded in length    & no  & yes & no  & yes & n/a & yes \\
constant per-frame cost       & no  & yes & no  & yes & n/a & yes \\
visible and thermal input     & yes & no  & no  & no  & yes & yes \\
small-object benchmark        & no  & no  & no  & no  & partly & yes \\
robust to misalignment        & no  & n/a & n/a & n/a & yes & yes \\
robust to sensor dropout      & no  & n/a & n/a & n/a & no  & yes \\
handles irregular frame rate  & no  & yes & no  & no  & n/a & yes \\
\end{tabularx}
\end{table}

\section{What the contribution is not}

Three claims will be avoided.
\begin{itemize}
  \item The work will not claim a new static fusion block as its
        contribution. The still-image \method{} is a baseline and an
        ablation, not the headline.
  \item It will not claim a new scan order.
  \item It will not claim state-of-the-art accuracy on COCO or on saturated
        fusion benchmarks.
\end{itemize}
The claim is narrower and stronger: a causal, constant-memory,
reliability-aware state-space memory for visible--thermal streams, together
with a robustness protocol that tests it. The searches behind this chapter
found no such method. That is evidence, not proof, and the search must be
repeated immediately before submission (Section~\ref{sec:submission}).
